\clearpage\thispagestyle{otFront}\phantomsection\label{mode-emploi}\pdfbookmark[0]{Mode d’emploi}{bm-mode-emploi}\otPageTitre{Mode d’emploi}{Des repères stables dans tout le livre}\fontsize{11}{15}\selectfont

\begin{tabularx}{\linewidth}{@{}lY@{}}\toprule
Une pastille & Découvrir : appliquer une propriété, lire une représentation ou faire un calcul direct.\\[3mm]
Deux pastilles & S’entraîner : choisir une méthode, combiner deux étapes ou interpréter un résultat.\\[3mm]
Trois pastilles & Résoudre : organiser une démarche plus longue, justifier et contrôler.\\\bottomrule
\end{tabularx}\par\medskip
Ces niveaux sont des repères pédagogiques, pas un classement des élèves.
Un exercice simple peut révéler une notion à consolider ; un problème peut
être réussi avec une représentation bien choisie.\par\medskip
\begin{exemple}[Une correction utile]Pour résoudre $3(x-2)=x+4$, on développe : $3x-6=x+4$, puis $2x=10$, donc $x=5$. Le contrôle dans l’énoncé donne $3(5-2)=9=5+4$. L’égalité obtenue vérifie la solution.\end{exemple}
\begin{attention}Lire tout l’énoncé, y compris les unités, les mots « exactement », « au moins », « et », « ou » et la précision demandée. En géométrie, un angle droit dessiné sans hypothèse n’est pas une preuve.\end{attention}
Les liens du sommaire et les numéros de page conduisent aux chapitres et
aux corrigés. Le tableau de bord permet de noter ce qui est réussi seul,
réussi avec aide ou à reprendre. Pour les scripts, les blocs indentés
appartiennent au bloc qui les contient ; les lignes « fin » en indiquent la fermeture.

\clearpage
